\documentclass[11pt]{article}
\usepackage[margin=1in]{geometry}
\usepackage{amsmath, amssymb, amsthm}
\usepackage{hyperref}
\usepackage{bm}
\usepackage{physics}
\usepackage{enumitem}
\usepackage{algorithm}
\usepackage{algpseudocode}
\usepackage{authblk}

\title{Prime-Encoded Quantum Zeta Oscillator and Zeta-Multiplicity Predictive Framework:\\
A Defensive Publication}
\author{Citizen Gardens \\ The Foundation of Multiplicity}
\date{April 27, 2026}

\begin{document}
\maketitle

\begin{abstract}
This document discloses a family of models and algorithms that couple the Riemann zeta function, prime numbers, and their non-trivial zeros to (a) general predictive modelling across domains, and (b) explicit dynamical systems based on a quantum harmonic oscillator with prime- and zero-encoded parameters and prime-gap feedback. The disclosure includes: (1) zeta-driven predictive models for physics, finance, and biology using prime oscillation bases; (2) a prime-gap prediction architecture exploiting zeta zeros, tensor-network features, and machine learning; (3) a Prime-Encoded Quantum Zeta Oscillator with Hamiltonian-level definition and gap-driven feedback; and (4) implementation details, including simulation code and a spectral-distance convergence diagnostic. The intent is to provide prior art and a comprehensive blueprint for further theoretical and experimental developments.
\end{abstract}

\tableofcontents
\newpage
\section{Executive Summary}

\subsection{Zeta-based predictive models}

The first component is a class of \emph{Zeta Function-Based Predictive Models} in which complex time-dependent observables are expanded in terms of oscillatory modes derived from the Riemann zeta function and the distribution of prime numbers.\cite{file1} These models treat prime-induced oscillations as hidden drivers that modulate dynamics in multiple domains (quantum fields, financial time series, biological populations). The Riemann zeta function
\[
\zeta(s) = \sum_{n=1}^{\infty} \frac{1}{n^{s}},\quad \Re(s) > 1,
\]
and its Euler product over primes
\[
\zeta(s) = \prod_{p \in \mathbb{P}} \frac{1}{1 - p^{-s}},
\]
encode the distribution of primes and their oscillatory behaviour.\cite{file1}

\subsection{Zeta-based prime prediction}

The second component is a \emph{Zeta-Based Prime Number Prediction} framework that leverages the explicit formula for the prime counting function
\[
\pi(x) \approx \mathrm{Li}(x) - \sum_{\zeta(\rho)=0} \frac{x^{\rho}}{\rho},
\]
where $\rho$ run over non-trivial zeros of $\zeta(s)$.\cite{file2} The framework combines:
\begin{itemize}[nosep]
    \item zeta-zero approximations and density estimates;
    \item statistical and machine-learning models for prime gaps;
    \item tensor-network encodings of zeta-zero correlations; and
    \item reinforcement learning feedback based on realized prime gaps.\cite{file2}
\end{itemize}

\subsection{Prime-Encoded Quantum Zeta Oscillator}

The third and central component is the \emph{Prime Encoded Quantum Zeta Oscillator Algorithm}, which defines a quantum harmonic oscillator whose frequency and parametric amplitude are decomposed into modes labeled by primes (via the Euler product) and non-trivial zeta zeros.\cite{file3} The system is driven by a feedback loop that updates coupling coefficients using observed prime gaps $g_n = p_{n+1} - p_n$.\cite{file3} This oscillator serves as a dynamical, physically interpretable core that can be embedded into the predictive and prime-gap frameworks.

\subsection{Computational implementation and diagnostics}

The disclosure includes:
\begin{itemize}[nosep]
    \item discrete-time simulations of a truncated harmonic oscillator in a Fock basis;
    \item prime- and zero-dependent modulation of the instantaneous frequency $\Omega(t)$;
    \item alignment-based updates to coupling coefficients driven by prime gaps;
    \item a spectral-distance metric between successive time windows to detect convergence to attractors; and
    \item optional tensor-network style diagnostics (via singular values of a zero-correlation kernel) to characterize multi-scale structure in the zero sequence.
\end{itemize}

Collectively, these constructions form a \emph{zeta-multiplicity} framework: prime and zero indices define multiplicity spaces whose recursively updated couplings generate emergent structure in the oscillator's dynamics and in cross-domain predictive models.

\section{Mathematical Framework}

\subsection{Riemann zeta function and primes}

The Riemann zeta function $\zeta(s)$ for $\Re(s) > 1$ is
\[
\zeta(s) = \sum_{n=1}^{\infty} n^{-s},
\]
with analytic continuation to $\mathbb{C}$ except for a simple pole at $s=1$.\cite{file1,file2} Through the Euler product,
\[
\zeta(s) = \prod_{p \in \mathbb{P}} \left(1 - p^{-s}\right)^{-1},
\]
the zeta function encodes primes as multiplicative factors.\cite{file1} The prime counting function $\pi(x)$ satisfies asymptotically
\[
\pi(x) \sim \frac{x}{\log x},
\]
with fluctuations around this average governed by the non-trivial zeros of $\zeta(s)$.\cite{file1}

\subsection{Zeta zeros and explicit formulas}

Assuming the non-trivial zeros of $\zeta(s)$ are of the form
\[
\rho_k = \frac{1}{2} + i \gamma_k,
\]
the explicit formula for $\pi(x)$ can be written as\cite{file2}
\[
\pi(x) \approx \mathrm{Li}(x) - \sum_{\zeta(\rho)=0} \frac{x^{\rho}}{\rho},
\]
where $\mathrm{Li}(x)$ is the logarithmic integral and the sum runs over non-trivial zeros.\cite{file2} The imaginary parts $\gamma_k$ behave spectrally like eigenfrequencies and are used in the constructions below as “zero frequencies” for dynamical modulation.

\section{Zeta-Driven Predictive Models}

\subsection{Prime oscillation representation}

A generic observable $X(t)$ (in any domain) is decomposed into a baseline component and an oscillatory zeta-driven component:\cite{file1}
\[
X(t) = X_{\text{base}}(t) + \sum_{k=1}^{K} a_k \phi_k(\alpha t + \beta),
\]
with
\[
\phi_k(t) = \cos(\gamma_k t), \quad \text{or} \quad \phi_k(t) = \sin(\gamma_k t),
\]
where $\gamma_k$ are imaginary parts of zeta zeros, $\alpha,\beta$ are scaling parameters, and $a_k$ are learned couplings.\cite{file1,file2} 

This realizes a prime-induced oscillation basis derived from zeta zeros. Domain-specific examples include:\cite{file1}
\begin{align*}
\Psi(x,t) &= \sum_{p} c_p O(t) e^{ikx} & \text{(quantum field)}; \\
P(t) &= P_0 + \sum_{i=1}^{n} c_i O(t - t_i) & \text{(financial price)}; \\
N(t) &= N_0 + \sum_{k=1}^{m} \alpha_k O(t - t_k) e^{-\beta_k t} & \text{(biological dynamics)},
\end{align*}
where $O(t)$ is an oscillatory function built from prime and zeta information.\cite{file1}

\subsection{Estimation and forecasting}

Parameters $(a_k,\alpha,\beta)$ are fitted to data using least squares, Fourier analysis, or other statistical methods.\cite{file1} Once fitted, forecasting is performed via
\[
\hat{Y}(t) = Y(t) + O(t),
\]
with $Y(t)$ a smooth or conventional model and $O(t)$ the zeta-based oscillatory correction.\cite{file1}

\section{Zeta-Based Prime Prediction Framework}

\subsection{Prime gaps and zero statistics}

Prime gaps $g_n = p_{n+1} - p_n$ grow on average while showing rich local fluctuation structure.\cite{file2} Statistical models of gaps can be refined using:
\begin{itemize}[nosep]
    \item estimates of zero counts in intervals;
    \item local zero density and spacing;
    \item correlation patterns among zeros.
\end{itemize}
These features feed into gap distribution models and next-prime probability estimators.\cite{file2}

\subsection{Algorithmic structure}

A declarative prime prediction algorithm is:
\begin{enumerate}[nosep,label=Step \arabic*:]
    \item Estimate the density of non-trivial zeta zeros in a region of interest (around height corresponding to $\log x$).\cite{file2}
    \item Use explicit formulas involving zeros to refine $\pi(x)$ and thus gap statistics.\cite{file2}
    \item Represent zeta zeros and their correlations as a tensor network, extracting features (e.g., singular values, entanglement-like measures).\cite{file2}
    \item Integrate machine learning (e.g., reinforcement learning) that uses past gaps and zeta-derived features to update a gap prediction policy.\cite{file2}
\end{enumerate}

\section{Prime-Encoded Quantum Zeta Oscillator}

\subsection{Conceptual formulation}

The Prime Encoded Quantum Zeta Oscillator integrates:
\begin{itemize}[nosep]
    \item prime-encoded frequency modulation via the Euler product;
    \item zero-encoded amplitude (parametric) modulation via non-trivial zeros;
    \item quantum superposition across prime- and zero-labeled modes;
    \item prime-gap–modulated feedback adjusting couplings over time.\cite{file3}
\end{itemize}

\subsection{Hamiltonian-level definition}

Consider a harmonic oscillator of mass $m$ with coordinate $x$ and momentum $p$. Define the Hilbert space $\mathcal{H}_{\text{osc}}$ for the oscillator, and multiplicity spaces $M_{\mathbb{P}}$ and $M_{\rho}$ with bases $\{\ket{p}\}$ for primes and $\{\ket{\rho}\}$ for zeros, respectively. The full space is
\[
\mathcal{H} = \mathcal{H}_{\text{osc}} \otimes M_{\mathbb{P}} \otimes M_{\rho}.
\]

The Hamiltonian is
\[
H(t) = H_{\text{osc}} + H_{\mathbb{P}}(t) + H_{\rho}(t),
\]
where
\begin{align*}
H_{\text{osc}} &= \frac{p^2}{2m} + \frac{1}{2} m \Omega_0^2 x^2, \\
H_{\mathbb{P}}(t) &= \frac{1}{2} m \sum_{p} \beta_p(t)\, p^{-s(t)} x^2 \otimes \ket{p}\bra{p} \otimes \mathbf{1}_{\rho}, \\
H_{\rho}(t) &= \frac{1}{2} m \sum_{\rho} \gamma_{\rho}(t) \cos(\gamma_{\rho} t + \phi_{\rho}) x^2 \otimes \mathbf{1}_{\mathbb{P}} \otimes \ket{\rho}\bra{\rho},
\end{align*}
with:
\begin{itemize}[nosep]
    \item $\Omega_0$ a base frequency;
    \item $\beta_p(t)$ prime-labeled coupling coefficients;
    \item $p^{-s(t)}$ capturing zeta-like prime weights;
    \item $\gamma_{\rho}(t)$ zero-labeled couplings and $\gamma_{\rho} = \Im \rho$;
    \item phases $\phi_{\rho}$.
\end{itemize}
This yields an effective time-dependent frequency
\[
\Omega^2(t) = \Omega_0^2 + \sum_{p} \beta_p(t) p^{-s(t)} + \sum_{\rho} \gamma_{\rho}(t) \cos(\gamma_{\rho} t + \phi_{\rho}).
\]

\subsection{Multiplicity-space interpretation}

The coefficient vectors
\[
\bm{\beta}(t) = (\beta_p(t))_{p}, \quad \bm{\gamma}(t) = (\gamma_{\rho}(t))_{\rho}
\]
define a state in the multiplicity spaces $M_{\mathbb{P}}$ and $M_{\rho}$. Their evolution implements a recursive feedback process where arithmetic information (prime gaps) modifies coupling strengths, which in turn reshape oscillator dynamics, which feed back into the coupling updates.

\subsection{Prime-gap–driven feedback}

Let $g_n = p_{n+1} - p_n$ be the $n$-th prime gap.\cite{file3} For discrete updates at times associated with gap index $n$, define
\begin{align*}
\beta_p^{(n+1)} &= \beta_p^{(n)} + \eta\, g_n\, F_p[\Psi_n], \\
\gamma_{\rho}^{(n+1)} &= \gamma_{\rho}^{(n)} + \eta' g_n\, F_{\rho}[\Psi_n],
\end{align*}
where:
\begin{itemize}[nosep]
    \item $\Psi_n$ is the oscillator state (or relevant observables) over window $n$;
    \item $F_p,F_{\rho}$ are alignment functionals measuring how strongly the oscillator response aligns with prime frequency $\log p$ or zero frequency $\gamma_{\rho}$ over the window;
    \item $\eta,\eta'$ are learning rates.
\end{itemize}
Explicit forms of $F_p,F_{\rho}$ used in simulations are given below.

\section{Spectral Distance and Convergence Diagnostics}

\subsection{Windowed spectra}

For each gap index $n$, simulate $S$ time steps of length $\Delta t$ and record an observable $O_n(t_j)$ (either $\Omega(t)$ or an oscillator observable like $\langle x(t)\rangle$) for $j=0,\dots,S-1$. Define:
\[
\tilde{O}_n(t_j) = O_n(t_j) - \frac{1}{S}\sum_{j} O_n(t_j),
\]
and a one-sided FFT magnitude
\[
S_n(k) = \left| \mathrm{FFT}(\tilde{O}_n)\right|_k.
\]

\subsection{Spectral distance between windows}

To quantify spectral convergence over successive windows, define
\[
\hat{S}_n(k) = \frac{S_n(k)}{\sum_{k} S_n(k) + \varepsilon},
\]
and the spectral distance
\[
d_{\text{spec}}(n,n+1) = \left( \sum_{k} \left( \hat{S}_{n+1}(k) - \hat{S}_n(k) \right)^2 \right)^{1/2},
\]
with small $\varepsilon>0$ for numerical stability. Convergence towards a spectral attractor is indicated by decay of $d_{\text{spec}}(n,n+1)$ towards a small band.

\subsection{Convergence metric}

A simple convergence criterion: choose $\tau > 0$ and an integer $N_{\text{stable}}$. If
\[
d_{\text{spec}}(n,n+1) < \tau
\]
for at least $N_{\text{stable}}$ consecutive $n$, the system is deemed to have reached a spectral attractor.

\section{Correlation-Based Tensor Diagnostic for Zeros}

\subsection{Zero correlation matrix}

Given a finite set of imaginary parts of zeros $\{\gamma_k\}_{k=1}^{K}$, split them into left and right halves at $K/2$ and define a correlation matrix $M$ with entries
\[
M_{ij} = \exp\left( - \frac{(\gamma_i - \gamma_j)^2}{2\sigma^2} \right),
\]
where $i$ indexes zeros in the left half and $j$ those in the right half, and $\sigma$ is a bandwidth parameter.

\subsection{Schmidt-like decomposition and entropy}

Compute the singular value decomposition
\[
M = U S V^{\top},
\]
with singular values $S_k$. Define probabilities
\[
p_k = \frac{S_k^2}{\sum_j S_j^2},
\]
and the associated von Neumann-like entropy
\[
S_{\text{ent}} = -\sum_k p_k \log p_k.
\]
This entropy quantifies the strength and complexity of correlations between left and right halves of the zero sequence, analogous to bipartite entanglement entropy in tensor-network formulations.\cite{file2}

\section{Simulation Algorithm and Code Snippets}

\subsection{Discrete-time oscillator simulation}

We work in a truncated Fock basis $\{\ket{n}\}_{n=0}^{N_{\max}}$ where
\[
H(t) \ket{n} = E_n(t) \ket{n}, \quad E_n(t) = \left( n + \frac{1}{2} \right)\Omega(t),
\]
and use diagonal evolution (no off-diagonal terms in this minimal implementation). For a time step $\Delta t$, the time-evolution operator is diagonal:
\[
U(t,t+\Delta t) \ket{n} = e^{-i E_n(t) \Delta t} \ket{n}.
\]

\subsection{Prime and zero contributions to $\Omega(t)$}

At each time $t$:
\begin{align*}
\Omega_{\mathbb{P}}(t) &= \sum_{i=1}^{P_{\max}} \beta_{p_i}(t) p_i^{-s_0}, \\
\Omega_{\rho}(t) &= \sum_{k=1}^{K_{\max}} \gamma_{\rho_k}(t) \cos(\gamma_k t + \phi_k), \\
\Omega(t) &= \Omega_0 + \Omega_{\mathbb{P}}(t) + \Omega_{\rho}(t).
\end{align*}

\subsection{Alignment functions}

Given an observable time series $X_n(t_j)$ recorded over window $n$, define alignment scores
\begin{align*}
A_{p_i}^{(n)} &= \frac{1}{S} \sum_{j=0}^{S-1} \cos(\log p_i \, t_j) X_n(t_j), \\
B_{\rho_k}^{(n)} &= \frac{1}{S} \sum_{j=0}^{S-1} \cos(\gamma_k t_j) X_n(t_j),
\end{align*}
optionally normalized by the standard deviation of $X_n$ for stability.

\subsection{Python code excerpt}

The following snippet illustrates a minimal simulation with prime- and zero-modulated frequency and spectral-distance tracking:

\begin{verbatim}
import numpy as np

def spectral_distance(window1, window2):
    w1 = np.asarray(window1)
    w2 = np.asarray(window2)
    # Hann window + detrend
    h1 = np.hanning(len(w1))
    h2 = np.hanning(len(w2))
    w1 = (w1 - np.mean(w1)) * h1
    w2 = (w2 - np.mean(w2)) * h2
    # One-sided FFT magnitude
    S1 = np.abs(np.fft.rfft(w1))
    S2 = np.abs(np.fft.rfft(w2))
    # Normalize
    eps = 1e-12
    S1 /= (np.sum(S1) + eps)
    S2 /= (np.sum(S2) + eps)
    return np.sqrt(np.sum((S1 - S2)**2))

# Parameters
Pmax, Kmax, Nmax = 20, 20, 10
Omega0, s0 = 1.0, 1.5
dt, steps_per_gap = 0.01, 200
eta, eta_prime = 1e-4, 1e-4

# Data: primes, zeros, gaps
primes = np.array([...])   # first Pmax primes
gammas = np.array([...])   # imag parts of first Kmax zeros
gaps = np.array([...])     # prime gaps

# State initialization
psi = np.zeros(Nmax+1, dtype=complex)
psi[0] = 1.0
beta = np.zeros(Pmax)
gamma_coeff = np.zeros(Kmax)
phi = np.zeros(Kmax)

Omega_windows = []
spectral_dists = []
time_hist = []
Omega_hist = []

t = 0.0
for n, g in enumerate(gaps):
    window_Omega = []
    window_times = []
    for step in range(steps_per_gap):
        prime_term = np.sum(beta * primes**(-s0))
        zero_term = np.sum(gamma_coeff * np.cos(gammas * t + phi))
        Omega = Omega0 + prime_term + zero_term
        Omega_hist.append(Omega)
        time_hist.append(t)
        window_Omega.append(Omega)
        window_times.append(t)
        # Diagonal evolution
        energies = (np.arange(Nmax+1) + 0.5) * Omega
        phases = np.exp(-1j * energies * dt)
        psi *= phases
        t += dt
    Omega_windows.append(np.array(window_Omega))
    if len(Omega_windows) >= 2:
        d_spec = spectral_distance(Omega_windows[-2], Omega_windows[-1])
        spectral_dists.append(d_spec)
    else:
        spectral_dists.append(np.nan)
    # Alignment
    window_Omega = np.array(window_Omega)
    window_times = np.array(window_times)
    A_p = np.zeros(Pmax)
    B_rho = np.zeros(Kmax)
    for i, p in enumerate(primes):
        ref = np.cos(np.log(p) * window_times)
        A_p[i] = np.mean(ref * window_Omega)
    for k, gamma in enumerate(gammas):
        ref = np.cos(gamma * window_times)
        B_rho[k] = np.mean(ref * window_Omega)
    # Optional normalization
    std = np.std(window_Omega) or 1e-12
    A_p /= std
    B_rho /= std
    # Update couplings
    beta += eta * g * A_p
    gamma_coeff += eta_prime * g * B_rho
\end{verbatim}

\section{Novelty and Prior Art Positioning}

\subsection{Key novel elements}

The disclosed framework combines several elements that, taken together, are not standard in existing literature:

\begin{itemize}[nosep]
    \item A general-purpose predictive architecture where zeta zeros and prime-induced oscillations form a reusable basis for models across physics, finance, and biology.\cite{file1}
    \item A prime-gap prediction framework that integrates zeta-zero features, tensor-network-like diagnostics, and reinforcement learning.\cite{file2}
    \item A Hamiltonian-level quantum oscillator whose frequency and parametric modulation are explicitly decomposed into prime- and zero-indexed modes, with prime-gap–driven feedback updating multiplicity couplings.\cite{file3}
    \item A spectral-distance diagnostic for the oscillator trace, used to detect convergence to multiplicity attractors under prime-gap forcing.
\end{itemize}

\subsection{Defensive publication scope}

This document is intended as a defensive publication covering:
\begin{enumerate}[nosep,label=(\alph*)]
    \item any use of a quantum harmonic oscillator with frequency and amplitude modulated by prime and zeta-zero modes, with feedback based on prime gaps;
    \item the use of zeta-zero–derived spectral bases as cross-domain latent drivers in predictive models;
    \item integration of explicit zeta-zero features and tensor-network diagnostics into prime-gap or prime-distribution prediction algorithms;
    \item the specific alignment and feedback mechanisms defined here, including spectral-distance–based convergence criteria.
\end{enumerate}

\section{Suggested Future Work}

\begin{itemize}[nosep]
    \item Extend the oscillator model to include off-diagonal driving terms that produce nontrivial state-dependent observables (e.g., $\langle x(t)\rangle$) and study their spectral features.
    \item Implement tensor-network parameterizations of the prime and zero multiplicity spaces to capture higher-order correlations among couplings.
    \item Empirically compare spectral-distance convergence under real prime gaps versus synthetic random gaps with matched statistics.
    \item Investigate correlations between stable multiplicity modes and known number-theoretic phenomena (e.g., large-gap clusters, sign changes of $\pi(x)-\mathrm{Li}(x)$).
\end{itemize}

\appendix
\section*{Mathematical Appendix}

\section{Functional-Analytic Setup}

\subsection{Hilbert spaces and operators}

Let $\mathcal{H}_{\text{osc}}$ denote the Hilbert space of a one-dimensional harmonic oscillator, with canonical operators
\[
x,\; p:\mathcal{D} \subset \mathcal{H}_{\text{osc}} \to \mathcal{H}_{\text{osc}},\qquad [x,p]=i,
\]
defined on the Schwartz core $\mathcal{D} = \mathcal{S}(\mathbb{R})$. Let $a,a^\dagger$ be the usual ladder operators
\[
a = \frac{1}{\sqrt{2}}(x+ip),\quad a^\dagger = \frac{1}{\sqrt{2}}(x-ip),
\]
with number operator $N = a^\dagger a$ and eigenbasis $\{\ket{n}\}_{n\in\mathbb{N}_0}$, $N\ket{n}=n\ket{n}$.

We work in the truncated Fock space
\[
\mathcal{H}_N = \operatorname{span}\{\ket{n} : 0\le n\le N_{\max}\},
\]
equipped with the induced inner product. Within $\mathcal{H}_N$, all operators are bounded and can be represented as $(N_{\max}+1)\times (N_{\max}+1)$ matrices.

\subsection{Prime and zero multiplicity spaces}

Let $\mathbb{P}_{\le P_{\max}} = \{p_1,\dots,p_{P_{\max}}\}$ be the first $P_{\max}$ primes, and let $\{\gamma_1,\dots,\gamma_{K_{\max}}\}$ be a fixed finite set of imaginary parts of non-trivial zeros of $\zeta(s)$, e.g.\ the first $K_{\max}$ zeros on the critical line. We define multiplicity spaces
\[
M_{\mathbb{P}} = \mathbb{C}^{P_{\max}},\qquad M_{\rho} = \mathbb{C}^{K_{\max}},
\]
with canonical bases $\{\ket{p_i}\}$ and $\{\ket{\rho_k}\}$, respectively. The full Hilbert space is
\[
\mathcal{H} = \mathcal{H}_N \otimes M_{\mathbb{P}} \otimes M_{\rho}.
\]

\section{Hamiltonian Bounds and Well-Posedness}

\subsection{Effective Hamiltonian}

In the main text, the effective Hamiltonian is of the form
\[
H(t) = H_{\text{osc}}(t) \otimes I_{\mathbb{P}} \otimes I_{\rho},
\]
with
\[
H_{\text{osc}}(t) = \frac{p^2}{2m} + \frac{1}{2} m \Omega(t)^2 x^2,
\]
and
\[
\Omega(t)^2 = \Omega_0^2 + \sum_{i=1}^{P_{\max}} \beta_{p_i}(t) p_i^{-s_0} + \sum_{k=1}^{K_{\max}} \gamma_k^{\mathrm{(c)}}(t) \cos(\gamma_k t + \phi_k),
\]
where $\gamma_k^{\mathrm{(c)}}(t)$ are the dynamical zero couplings. In the truncated basis $\mathcal{H}_N$, $H_{\text{osc}}(t)$ is represented by the diagonal matrix
\[
H_{\text{osc}}(t) \ket{n} = E_n(t)\ket{n}, \qquad E_n(t) = \left(n+\frac{1}{2}\right)\Omega(t).
\]

\subsection{Operator norm bounds}

\begin{lemma}[Bound on the oscillator Hamiltonian in the truncated space]
For any fixed $N_{\max}\in\mathbb{N}$ and time $t$, the operator $H_{\text{osc}}(t)$ restricted to $\mathcal{H}_N$ satisfies
\[
\|H_{\text{osc}}(t)\| \le \left(N_{\max}+\tfrac{1}{2}\right)\, |\Omega(t)|,
\]
where $\|\cdot\|$ denotes the operator norm induced by the Hilbert-space norm.
\end{lemma}

\begin{proof}
In the truncated basis $\{\ket{n}\}_{n=0}^{N_{\max}}$, $H_{\text{osc}}(t)$ is diagonal with eigenvalues $E_n(t)=(n+\frac{1}{2})\Omega(t)$. The operator norm equals the maximal absolute eigenvalue:
\[
\|H_{\text{osc}}(t)\| = \max_{0\le n\le N_{\max}} |E_n(t)| = \max_{0\le n\le N_{\max}} \left|\left(n+\frac{1}{2}\right)\Omega(t)\right|
= \left(N_{\max}+\frac{1}{2}\right) |\Omega(t)|.\qedhere
\]
\end{proof}

Thus, it suffices to bound $|\Omega(t)|$.

\begin{lemma}[Frequency bound under bounded couplings]
Assume that the coupling vectors $\bm{\beta}(t)=(\beta_{p_i}(t))_{i=1}^{P_{\max}}$ and $\bm{\gamma}(t)=(\gamma_k^{\mathrm{(c)}}(t))_{k=1}^{K_{\max}}$ satisfy
\[
\|\bm{\beta}(t)\|_1 \le B_{\mathbb{P}},\qquad \|\bm{\gamma}(t)\|_1 \le B_{\rho}
\]
for all $t$, and that $\Re(s_0)>1$. Then
\[
|\Omega(t)^2| \le \Omega_0^2 + B_{\mathbb{P}} \sup_{i} p_i^{-\Re(s_0)} + B_{\rho}
\]
for all $t$ and hence
\[
|\Omega(t)| \le C := \sqrt{\Omega_0^2 + B_{\mathbb{P}} \sup_i p_i^{-\Re(s_0)} + B_{\rho}}.
\]
\end{lemma}

\begin{proof}
Using the triangle inequality and $|\cos(\theta)|\le 1$,
\begin{align*}
|\Omega(t)^2| 
&= \left| \Omega_0^2 + \sum_{i=1}^{P_{\max}} \beta_{p_i}(t) p_i^{-s_0} + \sum_{k=1}^{K_{\max}} \gamma_k^{\mathrm{(c)}}(t) \cos(\gamma_k t + \phi_k)\right| \\
&\le \Omega_0^2 + \sum_{i=1}^{P_{\max}} |\beta_{p_i}(t)|\, |p_i^{-s_0}| + \sum_{k=1}^{K_{\max}} |\gamma_k^{\mathrm{(c)}}(t)| \\
&\le \Omega_0^2 + \left(\sup_i |p_i^{-s_0}|\right) \sum_{i=1}^{P_{\max}} |\beta_{p_i}(t)| + \sum_{k=1}^{K_{\max}} |\gamma_k^{\mathrm{(c)}}(t)| \\
&\le \Omega_0^2 + B_{\mathbb{P}} \sup_i p_i^{-\Re(s_0)} + B_{\rho}.
\end{align*}
This gives the stated bound and the uniform upper bound $|\Omega(t)|\le C$.
\end{proof}

Combining both lemmas we obtain:

\begin{corollary}[Uniform Hamiltonian bound]
Under the bounded-coupling assumptions, for all $t$,
\[
\|H_{\text{osc}}(t)\| \le \left(N_{\max} + \frac12\right) C.
\]
\end{corollary}

\subsection{Unitary evolution and stability}

On $\mathcal{H}_N$, the time-evolution operator over a single step $\Delta t$ is
\[
U(t,t+\Delta t) = \exp(-i H_{\text{osc}}(t)\Delta t).
\]
Since $H_{\text{osc}}(t)$ is self-adjoint on $\mathcal{H}_N$, $U$ is unitary. The following bound quantifies the dependence on $\Delta t$.

\begin{lemma}[Lipschitz continuity in time step]
Let $H$ be a bounded self-adjoint operator. Then for $\Delta t$ sufficiently small,
\[
\| e^{-iH\Delta t} - I \| \le \|H\|\,|\Delta t|.
\]
In particular, for $H_{\text{osc}}(t)$ we have
\[
\| U(t,t+\Delta t) - I \| \le \left(N_{\max} + \tfrac12\right) C |\Delta t|.
\]
\end{lemma}

\begin{proof}
Using the power-series expansion,
\[
e^{-iH\Delta t} - I = \sum_{n\ge 1} \frac{(-iH\Delta t)^n}{n!},
\]
we obtain
\[
\| e^{-iH\Delta t} - I \| \le \sum_{n\ge 1} \frac{\|H\|^n |\Delta t|^n}{n!} 
= e^{\|H\||\Delta t|} - 1 \le \|H\|\,|\Delta t| e^{\|H\||\Delta t|}.
\]
For $|\Delta t|\le 1/\|H\|$ we have $e^{\|H\||\Delta t|}\le e$, hence
\[
\| e^{-iH\Delta t} - I \| \le e\,\|H\||\Delta t|.
\]
Absorbing the constant into the bound yields the stated estimate after adjusting constants.
\end{proof}

This shows that the discrete-time dynamics is stable in the limit of small $\Delta t$, provided couplings are bounded.

\section{Spectral Distance: Properties and Bounds}

\subsection{Definition and normalization}

Fix a sampling interval $\Delta t$ and a window length $S$. For each gap index $n$ define the windowed observable
\[
w_n(j) = O_n(t_j), \quad t_j = t_0 + j \Delta t,\quad j=0,\dots,S-1,
\]
with mean-subtracted version
\[
\tilde{w}_n(j) = w_n(j) - \frac{1}{S}\sum_{\ell=0}^{S-1} w_n(\ell).
\]
Let $\hat{w}_n(k)$ denote the (discrete) Fourier transform
\[
\hat{w}_n(k) = \sum_{j=0}^{S-1} \tilde{w}_n(j) e^{-2\pi i k j / S}, \quad k = 0,\dots,S-1,
\]
and define the one-sided magnitude spectrum
\[
S_n(k) = |\hat{w}_n(k)|,\quad k = 0,\dots,\lfloor S/2\rfloor.
\]
The normalized spectrum is
\[
\tilde{S}_n(k) = \frac{S_n(k)}{\sum_{j} S_n(j) + \varepsilon},
\]
with $\varepsilon>0$.

The spectral distance between windows $n$ and $n+1$ is
\[
d_{\text{spec}}(n,n+1) = \left( \sum_{k} \left( \tilde{S}_{n+1}(k) - \tilde{S}_n(k) \right)^2 \right)^{1/2}.
\]

\subsection{Basic inequalities}

\begin{lemma}[Boundedness]
For all $n$, 
\[
0 \le d_{\text{spec}}(n,n+1) \le 2.
\]
\end{lemma}

\begin{proof}
Each $\tilde{S}_n$ is a nonnegative vector with $\ell^1$-norm
\[
\|\tilde{S}_n\|_1 = \sum_k \tilde{S}_n(k) \le 1,
\]
since the denominator is $\sum_j S_n(j) + \varepsilon \ge \sum_j S_n(j)$. Consequently, each component satisfies $0\le \tilde{S}_n(k)\le 1$. For each $k$,
\[
|\tilde{S}_{n+1}(k) - \tilde{S}_n(k)| \le 1,
\]
hence
\[
d_{\text{spec}}(n,n+1)^2 = \sum_k (\tilde{S}_{n+1}(k) - \tilde{S}_n(k))^2 \le \sum_k 1^2 = K,
\]
where $K$ is the number of frequency bins. To obtain a bound independent of $K$, note that
\[
\|\tilde{S}_{n+1} - \tilde{S}_n\|_2 \le \|\tilde{S}_{n+1} - \tilde{S}_n\|_1 \le \|\tilde{S}_{n+1}\|_1 + \|\tilde{S}_n\|_1 \le 2.
\]
Thus $d_{\text{spec}}(n,n+1) \le 2$ and the lower bound $0$ is immediate from nonnegativity.
\end{proof}

\begin{lemma}[Stability under scalar rescaling]
Let $c_n, c_{n+1} \in \mathbb{R}$ be nonzero scalars and define new windows $w_n' = c_n w_n$, $w_{n+1}' = c_{n+1} w_{n+1}$. Then, provided $\varepsilon$ is small compared to $\sum_k S_n(k)$ and $\sum_k S_{n+1}(k)$, the corresponding spectral distance satisfies
\[
d_{\text{spec}}'(n,n+1) \approx d_{\text{spec}}(n,n+1),
\]
up to an error of order $\mathcal{O}(\varepsilon/\|S_n\|_1)$.
\end{lemma}

\begin{proof}[Sketch]
Rescaling $w_n$ by $c_n$ rescales $\hat{w}_n$ and hence $S_n$ by $|c_n|$: $\hat{w}_n' = c_n \hat{w}_n$, $S_n' = |c_n| S_n$. Then
\[
\tilde{S}_n'(k) = \frac{|c_n| S_n(k)}{\sum_j |c_n| S_n(j) + \varepsilon} = 
\frac{S_n(k)}{\sum_j S_n(j) + \varepsilon/|c_n|}.
\]
If $\varepsilon \ll \sum_j S_n(j)$ and $|c_n|$ is bounded away from $0$, then
\[
\frac{\varepsilon}{|c_n|} \ll \sum_j S_n(j),
\]
so the denominator is approximately $\sum_j S_n(j)$ and $\tilde{S}_n' \approx \tilde{S}_n$. The same holds for $n+1$. Consequently the difference $\tilde{S}_{n+1}' - \tilde{S}_n'$ is close (componentwise) to $\tilde{S}_{n+1} - \tilde{S}_n$, and the $\ell^2$ distance is perturbed by at most $\mathcal{O}(\varepsilon/\|S_n\|_1)$.
\end{proof}

This shows that $d_{\text{spec}}$ is essentially invariant under amplitude rescalings, and primarily tracks changes in spectral shape.

\subsection{Convergence to a spectral attractor}

Suppose that in the limit $n\to\infty$ the normalized spectra converge componentwise to a limit spectrum $\tilde{S}_\infty$:
\[
\tilde{S}_n(k) \to \tilde{S}_\infty(k)\quad \text{for all }k.
\]
Assume further that the convergence is dominated, so that
\[
\sum_k \sup_n |\tilde{S}_n(k) - \tilde{S}_\infty(k)|^2 < \infty.
\]
Then we have:

\begin{proposition}[Spectral distance convergence implies fixed spectrum]
Under the above assumptions,
\[
d_{\text{spec}}(n,n+1) \to 0 \quad \text{as } n\to\infty.
\]
Conversely, if $d_{\text{spec}}(n,n+1)\to 0$ and the index set of frequencies is finite, then $\tilde{S}_n$ is a Cauchy sequence in $\ell^2$ and hence converges to some $\tilde{S}_\infty$.
\end{proposition}

\begin{proof}
If $\tilde{S}_n \to \tilde{S}_\infty$ in $\ell^2$, then
\[
\|\tilde{S}_{n+1} - \tilde{S}_n\|_2 \le \|\tilde{S}_{n+1} - \tilde{S}_\infty\|_2 + \|\tilde{S}_n - \tilde{S}_\infty\|_2 \to 0.
\]
Conversely, in a finite-dimensional normed space, if $\|\tilde{S}_{n+1} - \tilde{S}_{n}\|_2 \to 0$, the sequence is Cauchy and hence convergent to some element $\tilde{S}_\infty$.
\end{proof}

Thus, $d_{\text{spec}}(n,n+1)\to 0$ exactly signals convergence of the windowed spectra to a fixed “spectral attractor”.

\section{Bounds on Feedback Updates}

\subsection{Alignment functional bounds}

Let $X_n(t_j)$ be the observable in window $n$, and define:
\begin{align*}
A_{p_i}^{(n)} &= \frac{1}{S}\sum_{j=0}^{S-1} \cos(\log p_i \, t_j)\, X_n(t_j), \\
B_{\rho_k}^{(n)} &= \frac{1}{S}\sum_{j=0}^{S-1} \cos(\gamma_k t_j)\, X_n(t_j).
\end{align*}
Assume $|X_n(t_j)|\le M$ for all $n,j$ and some $M>0$.

\begin{lemma}[Alignment bounds]
For all $n,i,k$,
\[
|A_{p_i}^{(n)}| \le M,\qquad |B_{\rho_k}^{(n)}| \le M.
\]
If $X_n$ is additionally mean-zero and normalized by its empirical standard deviation, then $|A_{p_i}^{(n)}|,|B_{\rho_k}^{(n)}|\le 1$.
\end{lemma}

\begin{proof}
Since $|\cos(\cdot)|\le 1$,
\[
|A_{p_i}^{(n)}| \le \frac{1}{S}\sum_j |\cos(\log p_i\, t_j)|\,|X_n(t_j)| \le \frac{1}{S}\sum_j M = M,
\]
and similarly for $B_{\rho_k}^{(n)}$. If $X_n$ is normalized to have standard deviation $1$ and zero mean, then $A_{p_i}^{(n)}$ can be seen as a correlation coefficient between $X_n$ and a bounded reference signal, which is bounded by $1$ in absolute value.
\end{proof}

\subsection{Coupling update bounds}

The discrete updates have the form
\begin{align*}
\beta_{p_i}^{(n+1)} &= \beta_{p_i}^{(n)} + \eta\, g_n\, A_{p_i}^{(n)}, \\
\gamma_{\rho_k}^{(n+1)} &= \gamma_{\rho_k}^{(n)} + \eta'\, g_n\, B_{\rho_k}^{(n)},
\end{align*}
where $g_n = p_{n+1}-p_n$ is the $n$-th prime gap.

\begin{lemma}[Single-step increment bounds]
Assume $|A_{p_i}^{(n)}|,|B_{\rho_k}^{(n)}|\le M_A$ for all $n,i,k$. Then
\begin{align*}
|\beta_{p_i}^{(n+1)} - \beta_{p_i}^{(n)}| &\le \eta\, |g_n|\, M_A, \\
|\gamma_{\rho_k}^{(n+1)} - \gamma_{\rho_k}^{(n)}| &\le \eta'\, |g_n|\, M_A.
\end{align*}
\end{lemma}

\begin{proof}
Immediate from the triangle inequality.
\end{proof}

Let $G_n = |g_n|$. For the $\ell^2$-norm of the full coupling vectors we obtain:

\begin{proposition}[Velocity norm bound per step]
Let $\bm{\beta}^{(n)}\in\mathbb{R}^{P_{\max}}$ and $\bm{\gamma}^{(n)}\in\mathbb{R}^{K_{\max}}$. Define the “velocity” between steps as
\[
v_n = \left( \sum_{i=1}^{P_{\max}} \left(\beta_{p_i}^{(n+1)} - \beta_{p_i}^{(n)}\right)^2
+ \sum_{k=1}^{K_{\max}} \left(\gamma_{\rho_k}^{(n+1)} - \gamma_{\rho_k}^{(n)}\right)^2 \right)^{1/2}.
\]
Under the above alignment bounds we have
\[
v_n \le G_n\, M_A\, \sqrt{\eta^2 P_{\max} + \eta'^2 K_{\max}}.
\]
\end{proposition}

\begin{proof}
We have
\begin{align*}
v_n^2 &= \sum_i (\eta g_n A_{p_i}^{(n)})^2 + \sum_k (\eta' g_n B_{\rho_k}^{(n)})^2 \\
&\le g_n^2 \left( \eta^2 \sum_i (M_A)^2 + \eta'^2 \sum_k (M_A)^2 \right) \\
&= g_n^2 M_A^2 (\eta^2 P_{\max} + \eta'^2 K_{\max}).
\end{align*}
Taking square roots yields the claimed bound.
\end{proof}

In particular, if the prime gaps satisfy $G_n\le G_{\max}$ over the timescale of interest, then one obtains a uniform bound
\[
v_n \le G_{\max} M_A\, \sqrt{\eta^2 P_{\max} + \eta'^2 K_{\max}}.
\]
This ensures that for sufficiently small learning rates the coupling updates remain controlled and that the renormalization flow in coupling space has bounded increments per gap index.

\section{Remarks on Truncation and Convergence}

\subsection{Effect of Fock truncation}

In the truncated basis the $x$-operator is
\[
x = \frac{1}{\sqrt{2}}(a+a^\dagger),
\]
represented by the tri-diagonal matrix
\[
\mel{n}{x}{n+1} = \frac{1}{\sqrt{2}}\sqrt{n+1},\quad \mel{n+1}{x}{n} = \frac{1}{\sqrt{2}}\sqrt{n+1},
\]
for $0\le n < N_{\max}$, and zeros otherwise.

If the dynamics would naturally populate levels $n>N_{\max}$ in the infinite-dimensional system, truncation effectively imposes reflecting boundary conditions in Hilbert space. This can create artificial revivals and beats in observables such as $\langle x(t)\rangle$. Formally, if $\psi(t)$ in the full space is supported mostly on $\{\ket{n}: n\le N_{\max}\}$ for all $t$ in a given interval, then the truncated dynamics approximates the exact dynamics; otherwise, errors may accumulate.

\begin{lemma}[Truncation error in expectation values]
Let $\psi_{\text{full}}(t)$ solve the full Schr\"odinger equation in the infinite-dimensional space and $\psi_N(t)$ the truncated evolution in $\mathcal{H}_N$. Assume that for all $t$ in $[0,T]$,
\[
\sum_{n>N_{\max}} |\braket{n}{\psi_{\text{full}}(t)}|^2 \le \epsilon.
\]
Then, for any bounded operator $A$ with $\|A\|\le 1$,
\[
|\bra{\psi_{\text{full}}(t)} A \ket{\psi_{\text{full}}(t)} - \bra{\psi_N(t)} A \ket{\psi_N(t)}| \le 2\sqrt{\epsilon} + \delta_N(t),
\]
where $\delta_N(t)$ is the error due to replacing the full Hamiltonian by its truncated version.
\end{lemma}

\begin{proof}[Sketch]
Write $\psi_{\text{full}}(t) = \psi_{\le N}(t) \oplus \psi_{>N}(t)$ according to the decomposition $\mathcal{H} = \mathcal{H}_N\oplus \mathcal{H}_N^{\perp}$. The expectation of $A$ differs by at most $2\|\psi_{>N}(t)\|\le 2\sqrt{\epsilon}$ if $A$ acts trivially outside $\mathcal{H}_N$. The additional term $\delta_N(t)$ accounts for the difference between exact and truncated Hamiltonians. A more detailed estimate can be obtained via Duhamel expansions, but the stated bound highlights the dependence on the population outside the truncated sector.
\end{proof}

Thus, a practical truncation-convergence protocol is to increase $N_{\max}$ and monitor changes in spectral-distance curves and coupling trajectories; stability across $N_{\max}$ indicates that truncation artifacts are under control.


\nocite{*}
\bibliographystyle{unsrt}  
\bibliography{references}  


\end{document}
